\chapter{1987年全国卷（理）}

\section{单选题}

\begin{problem}

设 \(S\)，\(T\) 是两个非空集合，且 \(S \not\subset T\)，\(T \not\subset S\). 令 \(X = S \cap T\)，那么 \(S \cup X\) 等于（\quad）

\choices
    {\(X\)}
    {\(T\)}
    {\(\varnothing\)}
    {\(S\)}
\end{problem}

\begin{answer}
D
\end{answer}

\begin{solution}
由交集定义，\(X=S\cap T\) 的每个元素都属于 \(S\)，所以 \(X\subseteq S\). 因而由集合的吸收律，\(S\cup X=S\)，故选 D.
\end{solution}

\begin{problem}

设椭圆方程为 \(\frac{x^2}{a^2} + \frac{y^2}{b^2} = 1\)（\(a > b > 0\)），令 \(c = \sqrt{a^2 - b^2}\)，那么它的准线方程为（\quad）

\choices
    {\(y = \pm \frac{a^2}{c}\)}
    {\(y = \pm \frac{b^2}{c}\)}
    {\(x = \pm \frac{a^2}{c}\)}
    {\(x = \pm \frac{b^2}{c}\)}
\end{problem}

\begin{answer}
C
\end{answer}

\begin{solution}
椭圆的焦点为 \(\left(\pm c,0\right)\)，离心率为 \(e=\dfrac{c}{a}\). 由于长轴在 \(x\) 轴上，准线垂直于长轴，所以准线是两条竖直直线；每条准线到中心的距离为
\(\dfrac{a}{e}=\dfrac{a}{c/a}=\dfrac{a^2}{c}\). 因此准线方程为 \(x=\pm\dfrac{a^2}{c}\)，故选 C.
\end{solution}

\begin{problem}

设 \(a\)，\(b\) 是满足 \(ab < 0\) 的实数，那么（\quad）

\choices
    {\(|a + b| > |a - b|\)}
    {\(|a + b| < |a - b|\)}
    {\(|a - b| < ||a| - |b||\)}
    {\(|a - b| < |a| + |b|\)}
\end{problem}

\begin{answer}
B
\end{answer}

\begin{solution}
由 \(ab<0\) 可知 \(a\)、\(b\) 异号. 计算两边平方之差，得
\((a+b)^2-(a-b)^2=4ab<0\). 由于两个平方根均非负，所以 \(|a+b|<|a-b|\)，第二个选项正确. 又因为 \(a\)、\(b\) 异号，\(|a-b|=|a|+|b|\)，故选项 D 的严格不等式不成立；反三角不等式给出 \(|a-b|\geqslant\bigl||a|-|b|\bigr|\)，故选项 C 也不成立，选项 A 与第二个选项矛盾.
\end{solution}

\begin{problem}

已知 \(E\)，\(F\)，\(G\)，\(H\) 是空间中的四个点. 设命题甲：点 \(E\)，\(F\)，\(G\)，\(H\) 不共面；命题乙：直线 \(EF\) 和 \(GH\) 不相交. 那么（\quad）

\choices
    {甲是乙的充分条件，但不是必要条件}
    {甲是乙的必要条件，但不是充分条件}
    {甲是乙的充要条件}
    {甲不是乙的充分条件，也不是乙的必要条件}
\end{problem}

\begin{answer}
A
\end{answer}

\begin{solution}
若甲成立而直线 \(EF\) 与 \(GH\) 相交，则两条相交直线确定一个平面；\(E\)，\(F\) 在直线 \(EF\) 上，\(G\)，\(H\) 在直线 \(GH\) 上，于是四点都在这个平面内，与甲矛盾. 因此甲能推出乙，甲是乙的充分条件.

反之，取同一平面内的两条平行直线，令 \(E\)，\(F\) 在其中一条直线上，\(G\)，\(H\) 在另一条直线上，则四点共面，但 \(EF\) 与 \(GH\) 不相交. 所以乙不能推出甲，甲不是乙的必要条件. 故选 A.
\end{solution}

\begin{problem}

在区间 \((-\infty,0)\) 上为增函数的是（\quad）

\choices
    {\(y = -\log_{\frac{1}{2}}\left(-x\right)\)}
    {\(y = \frac{x}{1 - x}\)}
    {\(y = -(x + 1)^2\)}
    {\(y = 1 + x^2\)}
\end{problem}

\begin{answer}
B
\end{answer}

\begin{solution}
在 \((-\infty,0)\) 上，选项 A 可化为
\(y=-\log_{\frac12}\left(-x\right)=\dfrac{\ln\left(-x\right)}{\ln2}\)，所以 \(y'=\dfrac{1}{x\ln2}<0\)，不是增函数.

选项 B 中
\(y'=\dfrac{(1-x)+x}{(1-x)^2}=\dfrac{1}{(1-x)^2}>0\)，所以它在整个 \((-\infty,0)\) 上为增函数. 选项 C 的导数为 \(-2(x+1)\)，在 \(x=-1\) 两侧变号，不能在整个区间上递增；选项 D 的导数为 \(2x<0\)，为减函数. 故选 B.
\end{solution}

\begin{problem}

要得到函数 \(y = \sin \left(2x - \frac{\pi}{3}\right)\) 的图象，只需将函数 \(y = \sin 2x\) 的图象（如图）（\quad）

\choices
    {向左平行移动 \(\frac{\pi}{3}\)}
    {向右平行移动 \(\frac{\pi}{3}\)}
    {向左平行移动 \(\frac{\pi}{6}\)}
    {向右平行移动 \(\frac{\pi}{6}\)}
\end{problem}

\begin{answer}
D
\end{answer}

\begin{solution}
将自变量写成平移后的形式：
\(y=\sin\left(2x-\dfrac{\pi}{3}\right)=\sin\left(2\left(x-\dfrac{\pi}{6}\right)\right)\). 若 \(f(x)=\sin2x\)，则新函数为 \(f\left(x-\dfrac{\pi}{6}\right)\)，其图象由 \(y=f(x)\) 向右平移 \(\dfrac{\pi}{6}\) 得到. 故选 D.
\end{solution}

\begin{problem}

极坐标方程 \(\rho = \sin \theta + 2\cos \theta\) 所表示的曲线是（\quad）

\choices
    {直线}
    {圆}
    {双曲线}
    {抛物线}
\end{problem}

\begin{answer}
B
\end{answer}

\begin{solution}
由极坐标与直角坐标的关系 \(x=\rho\cos\theta\)、\(y=\rho\sin\theta\)，原方程两边乘以 \(\rho\)，得
\(\rho^2=\rho\sin\theta+2\rho\cos\theta=y+2x\). 又有 \(\rho^2=x^2+y^2\)，所以
\(x^2+y^2=y+2x\)，即
\(\left(x-1\right)^2+\left(y-\dfrac12\right)^2=\dfrac54\).
这是圆心为 \(\left(1,\dfrac12\right)\)、半径为 \(\dfrac{\sqrt5}{2}\) 的圆，故选 B.
\end{solution}

\begin{problem}

函数 \(y = \arccos(\cos x)\)（\(x \in \left[-\frac{\pi}{2}, \frac{\pi}{2}\right]\)）的图象是（\quad）

\choices
    {\choicebitmap[width=0.15\paperwidth]{img/1987/national_science/q8_choiceA.png}}
    {\choicebitmap[width=0.15\paperwidth]{img/1987/national_science/q8_choiceB.png}}
    {\choicebitmap[width=0.15\paperwidth]{img/1987/national_science/q8_choiceC.png}}
    {\choicebitmap[width=0.15\paperwidth]{img/1987/national_science/q8_choiceD.png}}
\end{problem}

\begin{answer}
A
\end{answer}

\begin{solution}
在 \(\left[-\dfrac{\pi}{2},\dfrac{\pi}{2}\right]\) 上，\(\cos x\) 的值属于 \([0,1]\)，且反余弦的值域为 \([0,\pi]\). 对 \(x\geqslant0\)，有 \(\arccos(\cos x)=x\)；对 \(x\leqslant0\)，有 \(\arccos(\cos x)=-x\). 因而
\(y=\left|x\right|\)，即图象由 \(y=-x\)（\(-\dfrac{\pi}{2}\leq x\leqslant0\)）和 \(y=x\)（\(0\leq x\leq\dfrac{\pi}{2}\)）组成. 它是顶点在原点、两端点纵坐标均为 \(\dfrac{\pi}{2}\) 的折线，故选 A.
\end{solution}

\section{解答题}

\begin{problem}

求函数 \(y = \tan \frac{2x}{3}\) 的周期.
\end{problem}

\begin{solution}
设 \(T>0\) 是该函数的最小正周期. 正切函数的最小正周期为 \(\pi\)，因此要求自变量 \(x\) 增加 \(T\) 后，正切函数的自变量增加 \(\pi\)，即
\(\dfrac{2(x+T)}{3}-\dfrac{2x}{3}=\dfrac{2T}{3}=\pi\). 解得 \(T=\dfrac{3\pi}{2}\)，故函数的最小正周期为 \(\dfrac{3\pi}{2}\).
\end{solution}

\begin{problem}

已知方程 \(\frac{x^2}{2 + \lambda} - \frac{y^2}{1 + \lambda} = 1\) 表示双曲线，求 \(\lambda\) 的范围.
\end{problem}

\begin{solution}
首先分母不能为零，所以 \(\lambda\ne-2\)，\(-1\). 当 \(2+\lambda\) 与 \(1+\lambda\) 同号时，方程左边的两个二次项符号相反，移项后可化为一个正项减去一个正项等于 \(1\)，表示双曲线. 当 \(2+\lambda>0\)、\(1+\lambda<0\) 时，两个二次项都为正，方程表示椭圆；当 \(2+\lambda<0\)、\(1+\lambda>0\) 时，左边两项都为负，不可能等于 \(1\).

所以表示双曲线的充要条件为
\((\lambda+2)(\lambda+1)>0\). 解此不等式，得 \(\lambda<-2\) 或 \(\lambda>-1\).
\end{solution}

\begin{problem}

若 \(\left(1 + x\right)^n\) 的展开式中，\(x^3\) 的系数等于 \(x\) 的系数的 7 倍，求 \(n\).
\end{problem}

\begin{solution}
按二项式展开的题意，\(n\) 取正整数. 若 \(n=1\) 或 \(n=2\)，则 \(x^3\) 的系数为 \(0\)，而 \(x\) 的系数为 \(n>0\)，不满足题设，所以 \(n\geqslant3\). 由二项式定理，\(x\) 的系数为 \(\mathrm C_n^1=n\)，\(x^3\) 的系数为 \(\mathrm C_n^3\). 因而
\(\mathrm C_n^3=7\mathrm C_n^1\)，即 \(\dfrac{n(n-1)(n-2)}{6}=7n\).

因为 \(n>0\)，可约去 \(n\)，得到 \((n-1)(n-2)=42\)，即 \(n^2-3n-40=0\). 解得 \(n=8\) 或 \(n=-5\)，只有正整数解 \(n=8\).
\end{solution}

\begin{problem}

求极限：\(\displaystyle\lim_{n \to \infty}\left(\frac{1}{n^2 + 1} + \frac{2}{n^2 + 1} + \frac{3}{n^2 + 1} + \cdots + \frac{2n}{n^2 + 1}\right)\).
\end{problem}

\begin{solution}
由于每一项的分母相同，先合并分子：
\(\dfrac{1+2+\cdots+2n}{n^2+1}=\dfrac{\dfrac{2n(2n+1)}{2}}{n^2+1}=\dfrac{n(2n+1)}{n^2+1}\). 分子、分母同除以 \(n^2\)，得到
\(\dfrac{2+\frac1n}{1+\frac1{n^2}}\). 当 \(n\to\infty\) 时，\(\dfrac1n\to0\)、\(\dfrac1{n^2}\to0\)，故原极限为
\(\dfrac{2+0}{1+0}=2\).
\end{solution}

\begin{problem}

在抛物线 \(y = 4x^2\) 上求一点，使该点到直线 \(y = 4x - 5\) 的距离为最短.
\end{problem}

\begin{solution}
设抛物线上动点为 \(P=(x,4x^2)\). 直线 \(y=4x-5\) 化为 \(4x-y-5=0\)，故
\[
d(P,l)=\frac{|4x-4x^2-5|}{\sqrt{17}}
=\frac{4x^2-4x+5}{\sqrt{17}}
=\frac{4\left(x-\frac12\right)^2+4}{\sqrt{17}}
\]
这里 \(4x^2-4x+5=4\left(x-\dfrac12\right)^2+4>0\)，所以绝对值确实可以按上式去掉. 由平方项的非负性，距离在且仅在 \(x=\dfrac12\) 时取得最小值；此时 \(y=4x^2=1\). 因此所求点为 \(\left(\dfrac12,1\right)\).
\end{solution}

\begin{problem}

由数字 1，2，3，4，5 组成没有重复数字且数字 1 与 2 不相邻的五位数，求这种五位数的个数.
\end{problem}

\begin{solution}
先不限制 \(1\) 和 \(2\) 是否相邻，五个互不相同的数字全排列共有 \(5!=120\) 个. 若 \(1\) 和 \(2\) 相邻，就把它们看作一个整体；这个整体内部有 \(12\)、\(21\) 两种次序，整体与数字 \(3\)，\(4\)，\(5\) 共四个对象，排列数为 \(4!\). 所以相邻的五位数有 \(2\cdot4!=48\) 个.

“不相邻”是“总数”去掉“相邻”，故所求五位数共有 \(5!-2\cdot4!=120-48=72\) 个.
\end{solution}

\begin{problem}

一个正三棱台的下底和上底的周长分别为 \(30 \mathrm{~cm}\) 和 \(12 \mathrm{~cm}\)，而侧面积等于两底面积之差，求斜高.
\end{problem}

\begin{solution}
正三棱台的下底、上底都是正三角形，故它们的边长分别为
\(\dfrac{30}{3}=10\mathrm{~cm}\)、\(\dfrac{12}{3}=4\mathrm{~cm}\). 正三角形面积为边长平方的 \(\dfrac{\sqrt3}{4}\) 倍，所以两底面积之差为
\(\dfrac{\sqrt3}{4}\left(10^2-4^2\right)=21\sqrt3\mathrm{~cm}^2\).
设斜高为 \(l\). 棱台的侧面积等于上下底周长和的一半乘斜高，故侧面积为
\(\dfrac{30+12}{2}l=21l\mathrm{~cm}^2\). 由题设 \(21l=21\sqrt3\)，得 \(l=\sqrt3\mathrm{~cm}\).
\end{solution}

\begin{problem}

求 \(\sin 10^\circ \sin 30^\circ \sin 50^\circ \sin 70^\circ\) 的值.
\end{problem}

\begin{solution}
利用恒等式
\[
4\sin\theta\sin\left(60^\circ+\theta\right)\sin\left(60^\circ-\theta\right)
=4\sin\theta\left(\sin^2 60^\circ-\sin^2\theta\right)
=3\sin\theta-4\sin^3\theta
=\sin3\theta
\]
恒等式中用到了 \(\sin(\alpha+\beta)\sin(\alpha-\beta)=\sin^2\alpha-\sin^2\beta\). 取 \(\theta=10^\circ\)，得
\(4\sin10^\circ\sin70^\circ\sin50^\circ=\sin30^\circ=\dfrac12\)，所以
\(\sin10^\circ\sin50^\circ\sin70^\circ=\dfrac18\). 再乘以 \(\sin30^\circ=\dfrac12\)，得到
\(\sin10^\circ\sin30^\circ\sin50^\circ\sin70^\circ=\dfrac12\cdot\dfrac18=\dfrac1{16}\).
\end{solution}

\begin{problem}

如图，三棱锥 \(P-ABC\) 中，已知 \(PA \perp BC\)，\(PA = BC = L\)，\(PA\)，\(BC\) 的公垂线 \(ED = h\). 求证：三棱锥 \(P-ABC\) 的体积 \(V = \frac{1}{6}L^2h\).
\end{problem}

\begin{solution}
由于 \(ED\) 是 \(PA\) 与 \(BC\) 的公垂线，故 \(ED\perp PA\)、\(ED\perp BC\)；再结合 \(PA\perp BC\)，这三条方向彼此垂直. 取 \(ED\)、\(PA\)、\(BC\) 的方向分别为 \(z\) 轴、\(x\) 轴、\(y\) 轴，并令
\(E=(0,0,0)\)，\(D=(0,0,h)\)，\(A=(p-L,0,0)\)，\(P=(p,0,0)\)，\(B=(0,q,h)\)，\(C=(0,q+L,h)\)，其中 \(p\)，\(q\) 为实数. 这些坐标满足 \(PA=BC=L\) 且 \(ED=h\).

由四面体的向量体积公式，
\(V=\frac16\left|\overrightarrow{PA}\cdot\left(\overrightarrow{PB}\times\overrightarrow{PC}\right)\right|\). 这里 \(\overrightarrow{PA}=(-L,0,0)\)，\(\overrightarrow{PB}=(-p,q,h)\)，\(\overrightarrow{PC}=(-p,q+L,h)\)，所以
\[
\overrightarrow{PB}\times\overrightarrow{PC}
=\left(qh-h(q+L),\ h(-p)-(-p)h,\ (-p)(q+L)-q(-p)\right)
=(-hL,0,-pL)
\]
因此
\(V=\frac16\left|(-L,0,0)\cdot(-hL,0,-pL)\right|=\frac16L^2h\)，命题得证.
\end{solution}

\begin{problem}

设对所有实数 \(x\)，不等式 \(x^2\log_2 \frac{4(a + 1)}{a} + 2x\log_2 \frac{2a}{a + 1} + \log_2 \frac{(a + 1)^2}{4a^2} > 0\) 恒成立，求 \(a\) 的取值范围.
\end{problem}

\begin{solution}
首先要求三个对数的真数均为正. 这三个条件等价于
\(\dfrac{a+1}{a}>0\)，所以 \(a>0\) 或 \(a<-1\)，并且 \(a\ne0\)，\(-1\). 令
\(r=\dfrac{2a}{a+1}\)，\(B=\log_2 r\). 在上述定义域内 \(r>0\)，且
\(\log_2\dfrac{4(a+1)}a=\log_2\dfrac8r=3-B\)，
\(\log_2\dfrac{(a+1)^2}{4a^2}=\log_2\dfrac1{r^2}=-2B\).

于是原不等式左边为关于 \(x\) 的二次式
\(f(x)=(3-B)x^2+2Bx-2B\). 要使 \(f(x)>0\) 对所有实数 \(x\) 恒成立，二次项系数必须为正，且判别式必须小于零；否则二次式在某处不正. 因而
\(3-B>0\)，且
\(\Delta=(2B)^2-4(3-B)(-2B)=4B(6-B)<0\).

第一条件给出 \(B<3\)，第二条件给出 \(B<0\) 或 \(B>6\)，合并得 \(B<0\). 因为底数 \(2>1\)，这等价于 \(0<r<1\).

若 \(a<-1\)，则
\(r-1=\dfrac{2a-(a+1)}{a+1}=\dfrac{a-1}{a+1}>0\)，与 \(r<1\) 矛盾. 若 \(a>0\)，则
\(\dfrac{2a}{a+1}<1\Longleftrightarrow 2a<a+1\Longleftrightarrow a<1\). 因而 \(a\) 的取值范围为 \(0<a<1\).
\end{solution}

\begin{problem}

设复数 \(z_1\) 和 \(z_2\) 满足关系式 \(z_1\overline{z_2} + \overline{A}z_1 + A\overline{z_2} = 0\)，其中 \(A\) 为不等于 \(0\) 的复数. 证明：

\begin{enumerate}
\item \(|z_1 + A||z_2 + A| = |A|^2\)；
\item \(\frac{z_1 + A}{z_2 + A} = \left|\frac{z_1 + A}{z_2 + A}\right|\).
\end{enumerate}
\end{problem}

\begin{solution}
\begin{enumerate}
\item 在已知等式两边加上 \(|A|^2=A\overline A\)，得到
\(z_1\overline{z_2}+\overline A z_1+A\overline{z_2}+|A|^2=|A|^2\). 左边正好因式分解为
\((z_1+A)\overline{(z_2+A)}=|A|^2\). 两边取绝对值，利用 \(|uv|=|u||v|\)，得
\(|z_1+A|\,|z_2+A|=|A|^2\). 由于 \(A\ne0\)，右边大于零，所以 \(z_1+A\ne0\)、\(z_2+A\ne0\).

\item 因为 \(z_2+A\ne0\)，可令 \(w=\dfrac{z_1+A}{z_2+A}\). 分子、分母同乘 \(\overline{z_2+A}\)，并利用上面的因式分解式，得
\(w=\dfrac{(z_1+A)\overline{(z_2+A)}}{|z_2+A|^2}=\dfrac{|A|^2}{|z_2+A|^2}>0\). 正实数的绝对值等于自身；也可由第（1）问直接算得
\(|w|=\dfrac{|z_1+A|}{|z_2+A|}=\dfrac{|A|^2}{|z_2+A|^2}=w\). 因而
\(\dfrac{z_1+A}{z_2+A}=w=|w|=\left|\dfrac{z_1+A}{z_2+A}\right|\).
\end{enumerate}
\end{solution}

\begin{problem}

设数列 \(a_1\)，\(a_2\)，\(\cdots\)，\(a_n\)，\(\cdots\) 的前 \(n\) 项的和 \(S_n\) 与 \(a_n\) 的关系是 \(S_n = -ba_n + 1 - \frac{1}{(1+b)^n}\)，其中 \(b\) 是与 \(n\) 无关的常数，且 \(b \neq -1\).

\begin{enumerate}
\item 求 \(a_n\) 和 \(a_{n-1}\) 的关系式；
\item 写出用 \(n\) 和 \(b\) 表示 \(a_n\) 的表达式；
\item 当 \(0 < b < 1\) 时，求极限 \(\displaystyle\lim_{n\to\infty} S_n\).
\end{enumerate}
\end{problem}

\begin{solution}
\begin{enumerate}
\item 由 \(S_n-S_{n-1}=a_n\)，并分别代入已知关系，得
\(a_n=-ba_n+ba_{n-1}+\dfrac{1}{(1+b)^{n-1}}-\dfrac{1}{(1+b)^n}\). 右端最后两项之差为
\(\dfrac{b}{(1+b)^n}\)，所以
\((1+b)a_n=ba_{n-1}+\dfrac{b}{(1+b)^n}\). 由于 \(b\ne-1\)，可以除以 \(1+b\)，得
\(a_n=\dfrac{b}{1+b}a_{n-1}+\dfrac{b}{(1+b)^{n+1}}\)，其中 \(n\geqslant2\).
\item 取 \(n=1\)，由 \(S_1=a_1\) 得
\(a_1=-ba_1+1-\dfrac{1}{1+b}=-ba_1+\dfrac{b}{1+b}\)，从而
\(a_1=\dfrac{b}{(1+b)^2}\). 令 \(u_n=(1+b)^{n+1}a_n\)，则第（1）问的关系式化为
\(u_n=bu_{n-1}+b\)，且 \(u_1=(1+b)^2a_1=b\). 若 \(u_k=b+b^2+\cdots+b^k\)，则
\(u_{k+1}=bu_k+b=b^2+b^3+\cdots+b^{k+1}+b=b+b^2+\cdots+b^{k+1}\). 结合 \(u_1=b\)，由数学归纳法得
\(u_n=b+b^2+\cdots+b^n\)，故
\(a_n=\dfrac{b+b^2+\cdots+b^n}{(1+b)^{n+1}}\).
\item 当 \(0<b<1\) 时，\(0<b+b^2+\cdots+b^n<\dfrac{b}{1-b}\)，而 \((1+b)^{n+1}\to+\infty\)，所以 \(a_n\to0\)；同时 \(\dfrac{1}{(1+b)^n}\to0\). 代回已知关系，得到
\(\displaystyle\lim_{n\to\infty}S_n=1\).
\end{enumerate}
\end{solution}

\begin{problem}

定长为 \(3\) 的线段 \(AB\) 的两端点在抛物线 \(y^2 = x\) 上移动，记线段 \(AB\) 的中点为 \(M\)，求点 \(M\) 到 \(y\) 轴的最短距离，并求此时点 \(M\) 的坐标.
\end{problem}

\begin{solution}
设 \(A=(u^2,u)\)、\(B=(v^2,v)\)，其中 \(u\)、\(v\) 为实数. 由 \(AB=3\)，得
\(9=(u^2-v^2)^2+(u-v)^2=(u-v)^2\left((u+v)^2+1\right)\). 记 \(s=u+v\)、\(d=u-v\)，则 \(d^2=\dfrac{9}{s^2+1}\). 中点 \(M\) 的横坐标为 \(x_M=\dfrac{u^2+v^2}{2}=\dfrac{s^2+d^2}{4}=\dfrac14\left(s^2+\dfrac{9}{s^2+1}\right)\).
因为 \(x_M\geqslant0\)，点 \(M\) 到 \(y\) 轴的距离就是 \(x_M\). 令 \(t=s^2\geqslant0\)，则
\(x_M=\dfrac14\left(t+\dfrac{9}{t+1}\right)\)，且
\(\dfrac{\mathrm d x_M}{\mathrm dt}=\dfrac14\left(1-\dfrac{9}{(t+1)^2}\right)\).
当 \(0\leqslant t<2\) 时导数为负，当 \(t>2\) 时导数为正，所以 \(t=2\) 给出全局最小值. 此时
\(x_M=\dfrac14\left(2+\dfrac93\right)=\dfrac54\).

又 \(s=\pm\sqrt2\)，且 \(d^2=\dfrac9{2+1}=3\)，故中点纵坐标
\(y_M=\dfrac{s}{2}=\pm\dfrac{\sqrt2}{2}\). 取 \(s=\sqrt2\)、\(d=\sqrt3\) 时，
\(u=\dfrac{s+d}{2}\)、\(v=\dfrac{s-d}{2}\) 均为实数，并且前面的长度等式给出 \(AB=3\)，所以该最小值确实可以取得. 因此最短距离为 \(\dfrac54\)，此时
\(M=\left(\dfrac54,\dfrac{\sqrt2}{2}\right)\) 或 \(M=\left(\dfrac54,-\dfrac{\sqrt2}{2}\right)\).
\end{solution}


\section{附加题}

\begin{problem}

求极限：\(\displaystyle\lim_{n\to\infty}\left(1 - \frac{1}{2x}\right)^x\).
\end{problem}

\begin{solution}
题面把极限下标印作 \(n\to\infty\)，而括号和指数中使用的变量为 \(x\). 按该表达式的含义，将趋于 \(+\infty\) 的变量统一记为 \(x\). 当 \(x\to+\infty\) 时，令 \(u=-\dfrac{1}{2x}\)，则 \(u\to0^-\)，且 \(x=-\dfrac{1}{2u}\). 因此
\[
\left(1-\frac{1}{2x}\right)^x
=(1+u)^{-\frac{1}{2u}}
=\left((1+u)^{\frac1u}\right)^{-\frac12}
\longrightarrow\e^{-\frac12}
\]
这里使用了基本极限 \(\displaystyle\lim_{u\to0^-}(1+u)^{1/u}=\e\)，且当 \(x\) 足够大时 \(1-\dfrac1{2x}>0\)，幂式有定义. 所以极限为 \(\e^{-\frac12}\).
\end{solution}

\begin{problem}

设 \(y = x\ln\left(1 + x^2\right)\)，求 \(y'\).
\end{problem}

\begin{solution}
由乘积求导法则和复合函数求导法则，
\[
y'=x'\ln\left(1+x^2\right)+x\left(\ln\left(1+x^2\right)\right)'
=\ln\left(1+x^2\right)+x\frac{2x}{1+x^2}
=\ln\left(1+x^2\right)+\frac{2x^2}{1+x^2}
\]
因此 \(y'=\ln\left(1+x^2\right)+\dfrac{2x^2}{1+x^2}\).
\end{solution}
